\documentclass{Phyport}

\usepackage{hyperref}
\hypersetup{
    colorlinks=true,
}
% 选用宏（如有额外需要的宏包请在此添加）
\usepackage{float}       % 提供 [H] 选项
\usepackage{booktabs}    % 三线表
\usepackage{amsmath}     % 数学公式
\usepackage{caption}     % 标题设置
\usepackage{siunitx}     % SI 单位排版
 %选用宏

\exname{} %实验名称
\extable{} %实验桌号
\instructor{} %指导教师
\class{电科2502} %班级
\name{管云哲} %姓名
\stuid{3250104519} %学号

\nyear{} %年
\nmonth{} %月
\nday{} %日
\nweekday{} %星期几

\redate{} %如有实验补做，补做日期
\resitu{} %情况说明：

\begin{document}
\setcounter{page}{1}
\makecover

\section{预习报告（10分）}
（注：将已经写好的“物理实验预习报告”内容拷贝过来）

\subsection{实验综述（5分）}
（自述实验现象、实验原理和实验方法，包括必要的光路图、电路图、公式等。不超过500字。）

\subsection{实验重点（3分）}
（简述本实验的学习重点，不超过100字。）

\subsection{实验难点（2分）}
（简述本实验的实现难点，不超过100字。）

\section{实验操作与原始数据（50分）}
（将有老师签名的“自备数据记录草稿纸”的扫描或手机拍摄图粘贴在下方，完整保留姓名，学号，教师签字和日期。按老师的要求拍摄实验过程图，附在此处）

\section{结果与分析（30分）}
\subsection{数据处理与结果（20分）}
（列出数据表格、选择适合的数据处理方法、写出测量或计算结果。）

\subsection{误差分析（10分）}
（运用测量误差、相对误差或不确定度等分析实验结果，写出完整的结果表达式，并分析误差原因。）

\section{实验探讨与思考题（10分）}
（对实验内容、现象和过程的小结，不超过100字。并解答教材或讲义或老师布置的思考题，请先写题干，再作答。）

{\fangsong 注意事项：
\begin{enumerate}[label=\arabic*.]
    \item 用PDF格式上传“实验报告”，文件名：学生姓名+学号+实验名称+周次。
    \item “实验报告”必须递交在“学在浙大”本课程内对应实验项目的“作业”模块内。
    \item “实验报告”成绩必须在“浙江大学物理实验教学中心网站”-“选课系统”内查询。
    \item 教学评价必须在“浙江大学物理实验教学中心网站”-“选课系统”内进行，学生必须进行教学评价，才能看到实验报告成绩，教学评价须在本次实验结束后3天内进行。
\end{enumerate}}
\begin{center}
    \fbfangsong{浙江大学物理实验教学中心制}
\end{center}
\end{document}
